\section{What India must fix}
\label{sec:india-fix}

Stated as work rather than as caveats, in descending order of how much they
cost.

\paragraph{1. No row of this study is both Indian and building-level.} The
Building Data Genome Project 2 harmonises $1{,}636$ buildings across three
countries and contains no Indian building, so the \emph{controller} was tested
on American buildings under an Indian tariff and the \emph{forecaster} on Indian
city demand, and no single row here is both. Relabelling an American building as
Indian would have been worse than the gap, so the gap is reported. An earlier
draft stated that no Indian building-level load dataset exists; that was wrong.
I-BLEND \citep{rashid2019iblend} provides 52 months (August 2013 to December
2017) of one-minute mains power and power factor for seven buildings on the
IIIT-Delhi campus --- academic, lecture, library, facilities, dining and two
dormitories --- under a CC0 licence. Three properties make it the right
corpus to close this gap and not merely a convenient one. It is in Delhi, the
city whose system demand is the study's focal series. It is sampled at one
minute, so the 15- and 30-minute billing blocks Indian demand charges are
assessed on can be formed by averaging rather than, as for every BDG2 building
here, interpolated up from hourly data that has already smoothed the peak. And
it spans four Junes, so the horizon-gap measurement of
Section~\ref{sec:horizon} can be replicated across years on one building rather
than reported once. What it lacks is submetering of HVAC, so the base-load
split remains a changepoint inference. The unchanged forecasting protocol has now been run on it
(Table~\ref{tab:iblend}); the closed loop under the Delhi (DERC) tariff has
not, and nothing in this paper anticipates its result.

\paragraph{2. Calibration collapses across the monsoon.} Coverage of $0.762$ is
an assumption failure, not a tuning failure: split conformal requires
calibration and test blocks to be exchangeable
\citep{vovk2005algorithmic,romano2019cqr}, and a Delhi May is not exchangeable
with a Delhi June. Section~\ref{sec:calibration} measures the remedy on this
feed: across a walk-forward year the adaptive layer, with its step carried over
from the building as a fraction of the interval width, holds $0.874$--$0.935$
by month where split conformal falls to $0.628$, and holds this June at
$0.895$; on the benchmark's own June, refitted and recalibrated
(Table~\ref{tab:aci-step}), it is $0.904$. The $0.762$ of
Table~\ref{tab:calibration} is the benchmark row under the step the building
used, and it is a units artefact rather than a monsoon one: the monsoon breaks
exchangeability, which split conformal cannot survive anywhere, and the layer
that survives it had been switched off by its units.

\paragraph{3. Indian load carries calendar structure the model does not encode.}
A large share of Indian public holidays are lunisolar --- Diwali, Holi, Eid ---
and move against the Gregorian calendar, so they cannot be expressed as the
fixed date list that load-forecasting pipelines conventionally use, including
ours until this study forced the change. Beyond holidays, Indian demand carries
load-shedding events, agricultural pumping schedules and festival periods that
have no counterpart in the European series and no representation in our feature
set. The low skill figure of $+0.097$ should be read partly as a statement about
missing features rather than solely about the data.

\paragraph{4. The Indian series is a decade old.} The Delhi archive covers
2011--2012. India's generation mix, air-conditioning penetration and peak demand
have all moved substantially since. The load \emph{shape} conclusions are the
durable ones; the levels are not, and nothing in this paper should be read as a
statement about the present Indian grid. We had a candidate answer to this
objection in the China arm and Section~\ref{sec:china} withdrew it.

\paragraph{What this does not undermine.} The favourable half of
Table~\ref{tab:india} rests on physical structure --- cooling share, the sign of
the load--temperature relationship, peakiness --- which is the part least likely
to have inverted since 2012. Indian demand has not stopped rising with
temperature.
